\section{Background: Divergence of Opinion}
\label{sec:background}

\epigraph{The very concept of uncertainty implies that reasonable men may
differ in their forecasts.}{---\citet{miller1977}}

\subsection{Miller's static premium hypothesis}
\label{sec:bg-miller}

\citet{miller1977} proposed a deceptively simple model of how heterogeneous
beliefs determine an asset's market price. Suppose any single investor can
purchase only one share---perhaps because of limited funds---and there are
$N$ shares available. The shares will end up owned by the $N$ investors with
the highest valuations. The marginal investor, the one with the lowest
valuation among the holders, sets the price. As the divergence of opinion
about the asset's value widens, this marginal valuation rises, pushing the
price above the mean estimate of all potential investors.

The implication is striking: when short selling is restricted, an asset's
price reflects the optimism of a minority, not the average view of the market.
Greater disagreement among investors thus produces a \emph{premium}, and
because eventual returns must be earned against this elevated price, future
returns on high-disagreement assets should be lower. This is the
``divergence of opinion premium hypothesis''---or, equivalently, the
\emph{overvaluation hypothesis} \citep{doukas2006}.

Miller offered initial empirical support: stocks about which there is the
greatest divergence of opinion at issuance experience smaller subsequent price
appreciation than ``seasoned'' stocks over horizons of one to five years.
Although Miller's model was largely ignored for decades, its narrative remains
unusually compelling. As \citet{moffitt2017} puts it, ``although Miller's
model has largely been ignored, it has lost none of its relevance.''

\subsection{The premium--discount debate}
\label{sec:bg-debate}

The empirical literature that grew up around Miller's hypothesis is
remarkable for its inability to converge. \citet{doukas2006} provide a careful
overview of the half-century of contested findings, of which the broad
outlines are as follows.

The \emph{discount} camp argues that divergence of opinion proxies for risk:
the higher the disagreement, the riskier the asset, and therefore the lower
its price relative to fundamentals \citep{williams1977,mayshar1983,
varian1985,merton1987,epsteinwang1994}. A direct consequence is that future
returns on high-disagreement assets should be \emph{higher}, not lower.

The \emph{premium} camp---supporting Miller's original prediction---finds
positive associations between dispersion in analyst earnings forecasts and
contemporaneous prices, with low subsequent returns. \citet{doukas2006}
themselves report evidence consistent with the discount hypothesis after
controlling for analyst-related dispersion measures, framing the apparent
contradictions as artefacts of measurement rather than substantive
disagreement.

Decades on, the debate is unresolved. \citet{moffitt2017} summarises the
state of play as follows:
\begin{quote}
``Though we have over two centuries of financial market history and various
theories of price formation, today there exists no single theory acceptable
to market practitioners, yet rigorous enough to satisfy market theorists.
The only major theory that purports universality is the efficient market
theory, which fails to explain some of the most important market phenomena,
e.g.\ bubbles and crashes.''
\end{quote}

\subsection{What is missing: dynamics}
\label{sec:bg-dynamics}

A common feature of Miller's model and its successors is that they are
fundamentally \emph{static}. They specify a one-shot relationship between the
distribution of beliefs and the market-clearing price. Real markets, however,
are continuous processes in which today's price feeds back into tomorrow's
beliefs, and in which the population of active buyers and sellers responds to
recent prices with adjustment lags.

This paper argues that the static framing is responsible for much of the
empirical confusion. When supply and demand are sample paths of a stochastic
process rather than fixed schedules, the same underlying mechanism---
heterogeneous valuations with random error---can produce \emph{either} a
divergence-of-opinion premium \emph{or} a discount, depending on the
balance between estimation noise and the market's adjustment capacity.
\Cref{sec:framework} formalises this claim.

\citet{shiller2003} reviews how behavioural finance has progressively
encroached on EMH territory; the present work occupies a related niche,
treating price formation itself as the locus of behavioural noise rather
than treating noise as a deviation from a frictionless benchmark.
